\documentclass[12pt]{ximera}
\title{Homework 9: Binomial Probability and the Subtraction Game}
\begin{document}
\begin{abstract}
Section 8.4 and the subtraction game: binomial probabilities for guessing on a
multiple-choice exam and for Steph Curry's three-point shooting, and a subtraction game
where the second player wins.
\end{abstract}
\maketitle

\section*{Section 8.4 and the Subtraction Game}

Throughout, the binomial probability formula is
\[
P(X=k) = C(n,k)\, p^k\, (1-p)^{\,n-k},
\]
where $n$ is the number of independent trials, $p$ is the probability of success on
any one trial, and $k$ is the number of successes.

\begin{problem}
A university's calculus exams consist of six multiple-choice questions, each with
five possible answers, followed by a few long-form questions.

\begin{prompt}
\textbf{(a)} If a student guesses randomly on the multiple-choice, what is their
probability of guessing correctly on each individual question?

As a fraction: $\answer{1/5}$

\begin{feedback}
One of the five answers is correct and the guess is random, so the probability is
$\frac15=0.2$.
\end{feedback}
\end{prompt}

\begin{prompt}
\textbf{(b)} What is the probability of guessing correctly on \emph{exactly four} of the
six multiple-choice questions? Give a percentage rounded to two decimal places.

$\answer[tolerance=0.005]{1.54}\%$

\begin{feedback}
With $n=6$, $p=0.2$, and $k=4$,
\[ P(X=4)=C(6,4)(0.2)^4(0.8)^2=15(0.0016)(0.64)=0.01536, \]
about $1.54\%$, or $2\%$ to the nearest whole percent.
\end{feedback}
\end{prompt}

\begin{prompt}
\textbf{(c)} What is the probability of guessing correctly on \emph{at least one} of the
six multiple-choice questions? Give a percentage rounded to two decimal places.

$\answer[tolerance=0.005]{73.79}\%$

\begin{hint}
``At least one'' is the complement of ``none''.
\end{hint}

\begin{feedback}
The only way to miss the event is to get all six wrong, which has probability
$(0.8)^6=0.262144$. So
\[ P(X\ge 1)=1-0.262144=0.737856, \]
about $73.79\%$, or $74\%$ to the nearest whole percent.
\end{feedback}
\end{prompt}
\end{problem}

\begin{problem}
Stephen Curry plays basketball for the Golden State Warriors and is well-known for
his exceptional shooting. Based on his 2024--25 regular season averages, he attempted
around eleven 3-point shots per game, and made about $40\%$ of the shots he took.

\begin{prompt}
\textbf{(a)} Complete Pascal's Triangle out to the eleventh row to solve the values
of $C(11,k)$ for every value of $k=0,1,2,\ldots,11$. Enter that row in order.

$\answer{1}$, $\answer{11}$, $\answer{55}$, $\answer{165}$, $\answer{330}$,
$\answer{462}$, $\answer{462}$, $\answer{330}$, $\answer{165}$, $\answer{55}$,
$\answer{11}$, $\answer{1}$

\begin{hint}
Each entry is the sum of the two entries above it, and every row starts and ends with
$1$. Row $10$ is $1, 10, 45, 120, 210, 252, 210, 120, 45, 10, 1$.
\end{hint}

\begin{feedback}
Row $11$ of Pascal's Triangle is
\[
1,\ 11,\ 55,\ 165,\ 330,\ 462,\ 462,\ 330,\ 165,\ 55,\ 11,\ 1.
\]
For example $C(11,4)=C(10,3)+C(10,4)=120+210=330$. Notice the row is symmetric,
since $C(11,k)=C(11,11-k)$.
\end{feedback}
\end{prompt}


\begin{prompt}
\textbf{(b)} What is the probability that he makes \emph{eight or more} out of eleven
attempted 3-point shots in a game?

As a percentage rounded to two decimal places: $\answer[tolerance=0.005]{2.93}\%$

\begin{hint}
``Eight or more'' means $k=8,9,10,11$. Use $C(11,8)=165$, $C(11,9)=55$,
$C(11,10)=11$, and $C(11,11)=1$ from part (a).
\end{hint}

\begin{feedback}
With $n=11$ and $p=0.4$:
\[
P(X=8)=165\,(0.4)^8(0.6)^3 \approx 0.023357,
\]
\[
P(X=9)=55\,(0.4)^9(0.6)^2 \approx 0.005190,
\]
\[
P(X=10)=11\,(0.4)^{10}(0.6)^1 \approx 0.000692,
\]
\[
P(X=11)=1\,(0.4)^{11} \approx 0.000042.
\]
Adding these gives
\[
P(X\ge 8)\approx 0.029281,
\]
about $2.93\%$, which rounds to $3\%$ at the nearest whole percent. A game like that
happens roughly two or three times per season.
\end{feedback}
\end{prompt}
\end{problem}

\begin{problem}
Alice and Bob are playing a version of the Subtraction Game where they start with a pile
of $54$ stones and on each turn they may take $1$, $2$, $3$, $4$, or $5$ stones. The
objective, as always, is to claim the last stone.

\begin{prompt}
\textbf{(a)} Describe the winning strategy in this version of the game.

\begin{freeResponse}
\end{freeResponse}

\begin{feedback}
Always leave your opponent a multiple of $6$. From a multiple of $6$, whatever your
opponent takes ($1$ to $5$), you take the rest of that group of $6$: if they take $k$,
you take $6-k$. The pile then drops by exactly $6$ each round, until your opponent faces
$0$ --- meaning you took the last stone.
\end{feedback}
\end{prompt}

\begin{prompt}
\textbf{(b)} If both players play optimally and Alice goes first, who will win?

\begin{multipleChoice}
\choice{Alice}
\choice[correct]{Bob}
\end{multipleChoice}

\begin{feedback}
Bob. The starting pile, $54=6\cdot 9$, is already a multiple of $6$, so Alice faces a
losing position. Whatever she takes, Bob answers with $6-k$ and keeps handing her a
multiple of $6$.
\end{feedback}
\end{prompt}
\end{problem}

\end{document}
